Mechanistic interpretability is increasingly used to guide interventions, such
as steering activations, removing circuits, or enforcing safety boundaries.
But an internal estimate that is accurate on average can still lead to a poor
decision when it is used to choose a specific action.

We present \textbf{ObserverBench}, a benchmark framework for testing whether
an internal estimator---an observer---is adequate for the intervention,
control, or safety task it directs. Each task fixes the model, information
available to the observer, allowed actions, decision rule, unseen test cases,
and loss. ObserverBench reports both estimation accuracy and downstream action
performance. Its rankings are task-specific: observers are compared only
under the same information and action constraints.

Our theory and experiments show three ways these criteria can differ:

\begin{itemize}[leftmargin=1.4em,itemsep=0.15em,topsep=0.25em]
  \item \textbf{Task-relevant accuracy.} In a control loop, observer errors
  matter at the starting point and along directions the allowed intervention
  can reach.

  \item \textbf{Prediction versus decision.} On circuit-intervention tasks in
  GPT-2-small and Qwen2.5-7B, pairwise observers predict unseen intervention
  effects more accurately without always choosing better actions. Observers
  trained to predict action loss directly choose lower-loss actions.

  \item \textbf{Consequence-aware safety.} Under fixed intervention budgets,
  a score that perfectly separates policy violations can allocate actions
  poorly when violations have different costs. Balanced safety panels can also
  conceal differences that appear when attacks are rare.
\end{itemize}

ObserverBench provides a common way to evaluate interpretability methods
through the actions they enable.
